\documentclass[conference]{IEEEtran}

% Preamble follows the IEEE conference template. Add only packages the content
% needs (below); do not add global spacing overrides to fit the page limit.
\usepackage{cite}
\usepackage{amsmath,amssymb,amsfonts}
\usepackage{graphicx}
\usepackage{textcomp}
\usepackage{xcolor}
\usepackage{url}        % URLs in references
\usepackage{booktabs}   % three-rule tables

\providecommand{\Description}[1]{}   % alt text (no-op in IEEEtran)
\pdfinfoomitdate=1                   % keep PDF metadata free of timestamps
\graphicspath{{figs/}}

\begin{document}

\title{<Title: phenomenon + setting + condition>}

\author{\IEEEauthorblockN{Anonymous submission}
\IEEEauthorblockA{Anonymous Institution\\Anonymous City, Anonymous Country}}

\maketitle

\begin{abstract}
% Five sentences: problem; finding with two numbers; mechanism; design response
% and its gate; boundary. Three memorable numbers, not ten.
\end{abstract}

\section{Introduction}
% Six paragraphs: assumption; finding + running example; why tuning keys cannot
% fix it; the different signal; evidence + boundary; contributions.

\section{Background and Model}

\section{Design}

\section{Implementation}

\section{Evaluation}
\paragraph{Setup}
% Paired unit, rounds analysed, floors, aggregation convention, interval type.

\section{Related Work}

\section{Conclusion}

\clearpage
\bibliographystyle{IEEEtran}
\bibliography{refs}

\end{document}
